\section{Genetic Algorithms and Search}
\label{sec:genetic_algorithms_pruning_only}

With the hardware-aware metrics introduced above, it becomes possible to compare model variants concretely. In particular, different designs may be compared with regard to prediction accuracy, parameter count, memory footprint, computational cost, memory bandwidth usage, and latency/throughput. In compression, these variants may correspond to different pruning masks, retained heads, feed-forward widths, removed layers, or layer-wise numerical precision.

Taken together, these choices form a large design space. Since enumerating all possible candidates is usually prohibitive, a search procedure is needed to explore this space and identify candidates that satisfy efficiency requirements while preserving prediction accuracy.

To manage this search, one can maintain a set of candidate designs instead of tracking a single design at a time. Each candidate corresponds to a possible compressed model. The candidates are evaluated, the more relevant ones are kept with higher probability, and new candidates are generated from existing ones through small changes or recombination. By repeating this process, the candidate set can move toward desirable regions of the design space.

A population-based search algorithm of this kind is known as a genetic algorithm. In this terminology, a candidate compressed model is an individual, the set of candidates is a population, and the evaluation measure is the fitness function. The following section introduces these notions together with encoding, selection, crossover, and mutation, then links them to compression design choices \cite{holland1975adaptation,goldberg1989genetic,deb2001multi}.

\subsection{Genetic Algorithm Fundamentals}

A Genetic Algorithm (GA) is a population-based search method inspired by biological evolution. It maintains candidate solutions, represents them in a form that can be modified, evaluates their quality, and repeatedly generates new candidates through selection, crossover, and mutation \cite{holland1975adaptation,goldberg1989genetic,michalewicz1996genetic}.

\subsubsection{Components of a Genetic Algorithm}

The basic object manipulated by a GA is an individual, also called a chromosome. Each individual encodes one candidate solution to the problem. A population is then a set of such individuals evaluated during the same generation. The fitness function assigns a quality score to each candidate, and the genetic operators define how new candidates are produced from existing ones \cite{goldberg1989genetic,michalewicz1996genetic}. These definitions establish what the algorithm manipulates; the next issue is how a candidate is written as a chromosome.

In model compression, an individual may represent a compressed architecture. For example, consider the chromosome
\begin{equation}
	\mathbf{c} = [1,0,1,1 \mid 2048,1024 \mid 8,4,4,8],
\end{equation}
where the first block indicates which attention heads are retained, the second block gives selected feed-forward widths, and the final block assigns layer-wise bit-widths. The fitness value then measures how good this candidate is with respect to the desired objective, such as validation accuracy, parameter count, memory footprint, or latency.

\subsubsection{Encoding Methods for Individuals}

The encoding determines how a candidate solution is written as a chromosome and therefore what kind of changes the algorithm can apply \cite{michalewicz1996genetic}. A \textbf{binary encoding} represents each gene with one bit, so every decision has two possible states. In compression, a bit value of $1$ may indicate that a neuron, head, or channel is retained, while $0$ indicates that it is removed. An \textbf{integer encoding} represents each gene with an integer selected from a finite set of choices. This is useful when a layer must choose one option among several possible bit-widths, such as $\{2,4,8,16\}$.

The choice of encoding is important because it determines what the genetic operators can change. A poor encoding may create invalid candidates or make useful solutions hard to reach. A good encoding reflects the structure of the compression problem while keeping the search space manageable. Once this representation is fixed, the algorithm can operate on a population across generations.

\subsubsection{Operation of Genetic Algorithms}

A GA begins by initializing a population, either randomly or from existing candidate designs. Each individual is evaluated using the fitness function. The algorithm then selects promising candidates, recombines their encoded decisions through crossover, and applies mutation to preserve exploration \cite{goldberg1989genetic}. The resulting candidates form the next population, and this cycle is repeated until a stopping criterion is reached, such as a maximum number of generations, convergence of the population, or an evaluation budget. The three operators below define how this update takes place.

\paragraph{Selection Operator.}
Selection chooses which individuals are allowed to reproduce. Better candidates should have a higher chance of being selected, but weaker candidates are not always discarded immediately because they may contain useful partial structures. Common strategies include roulette-wheel selection, rank-based selection, and tournament selection \cite{baker1985adaptive,miller1995genetic}. The selected candidates then serve as parents for crossover.

\paragraph{Crossover Operator.}
Crossover combines information from two or more parent individuals to produce offspring. In binary or integer encodings, this may exchange segments of chromosomes between parents. In uniform crossover, each gene can be inherited independently from either parent \cite{syswerda1989uniform}. For compression, crossover can combine useful pruning patterns or precision assignments discovered in different candidates. Since crossover reuses information already present in the population, mutation is used to introduce local changes.

\paragraph{Mutation Operator.}
Mutation introduces small random changes into an individual. It prevents the population from collapsing too early around one region of the search space and allows new structures to appear \cite{back1993optimal}. In a compression chromosome, mutation may flip a keep-or-remove bit, change the bit-width of a layer, or adjust a pruning threshold.

\vspace{0.4em}
\noindent\rule{0.35\textwidth}{0.4pt}
\vspace{0.4em}

With these mechanics defined, the compression problem can now be written as a genetic search problem.

\subsection{Compression as a Search Problem}

With the GA terminology in place, compression can be described as a search over encoded model variants. Let $\bx$ denote a compression configuration, such as a pruning mask, a vector of layer-wise bit-widths, or a set of structural choices. Evaluating $\bx$ produces a compressed model whose quality depends on both predictive performance and deployment cost.

In pruning, this means that a candidate can encode which Transformer structures remain active. For example, a chromosome may represent retained attention heads, selected FFN widths, or removed layers. The search procedure then compares complete candidates rather than isolated pruning decisions, which is important when the quality of one removal depends on the rest of the architecture.
